\documentclass[11pt]{article}
\usepackage[a4paper,margin=2.3cm]{geometry}
\usepackage{lmodern}
\usepackage[T1]{fontenc}
\usepackage{amsmath,amssymb,amsthm}
\usepackage{booktabs}
\usepackage{adjustbox}
\usepackage{graphicx}
\usepackage[hidelinks]{hyperref}

\renewcommand{\thesection}{S\arabic{section}}
\renewcommand{\thetable}{S\arabic{table}}
\renewcommand{\theequation}{S\arabic{equation}}

\theoremstyle{plain}
\newtheorem{theorem}{Theorem}
\newtheorem{lemma}[theorem]{Lemma}
\newtheorem{proposition}[theorem]{Proposition}
\theoremstyle{definition}
\newtheorem{definition}{Definition}

\newcommand{\R}{\mathbb{R}}
\newcommand{\E}{\mathbb{E}}
\newcommand{\Dhat}{\widehat{D}}
\newcommand{\TM}{\mathrm{TM}_b}
\newcommand{\DTM}{\Delta_{\mathrm{TM}}}
\newcommand{\gt}{\tilde g}
\newcommand{\rt}{CLAMP-RT}

\title{Supplementary material for\\``CLAMP-RT: Falsification-Robust Fairness Control for Federated Graph Learning via a Two-Server Trimmed Statistic with Attribute-Level Differential Privacy''}
\author{Quang-Vinh Dang and Ngoc-Son-An Nguyen}
\date{}

\setlength{\emergencystretch}{3em}
\begin{document}
\maketitle

\noindent Section S1 gives complete proofs of the results of Section~5 of the main text; theorem numbers follow the main text. Section S2 contains the full tables of the attack studies and the privacy sweep. All tables are generated by \texttt{make\_tables.py} from the result files listed in its header.

\section{Proofs}

\paragraph{Notation}
For $x\in\R^K$, $x_{(1)}\le\dots\le x_{(K)}$ are its sorted entries and $\TM(x)=\frac{1}{K-2b}\sum_{i=b+1}^{K-b}x_{(i)}$ for $0\le b$ with $K-2b\ge1$. For a validation set $\mathcal{V}_k$ split into groups $\mathcal{V}_{k,0},\mathcal{V}_{k,1}$ by the sensitive attribute, $m_{k,a}=\sum_{v\in\mathcal{V}_{k,a}}p(v)/\max(|\mathcal{V}_{k,a}|,n_{0,k})$ with $n_{0,k}=\max(1,\lceil c|\mathcal{V}_k|\rceil)$, and $d_k=\mathrm{clip}(m_{k,1}-m_{k,0},-1,1)$.

\subsection*{Lemma 1 (order statistics)}
\emph{Claim (a).} The map $x\mapsto x_{(i)}$ is non-decreasing in every entry and satisfies $(x+t\mathbf{1})_{(i)}=x_{(i)}+t$.

\emph{Proof.} We have $x_{(i)}=\max_{S\subseteq[K],|S|=K-i+1}\min_{l\in S}x_l$: for the set $S$ of indices of the $K-i+1$ largest entries the minimum is $x_{(i)}$, and any set of $K-i+1$ indices contains an index among the $i$ smallest entries, so its minimum is at most $x_{(i)}$. A maximum of minima of coordinates is non-decreasing in every coordinate and commutes with adding a constant to all coordinates. \qed

\emph{Claim (b).} If $x,y$ differ only in entry $j$, then all differences $y_{(i)}-x_{(i)}$ have the sign of $y_j-x_j$ (or vanish) and $\sum_i|y_{(i)}-x_{(i)}|=|y_j-x_j|$.

\emph{Proof.} Suppose $y_j\ge x_j$ (the other case is symmetric). Then $y\ge x$ entrywise, so $y_{(i)}\ge x_{(i)}$ for every $i$ by (a). Sorting is a permutation, so $\sum_i y_{(i)}=\sum_i y_i$ and $\sum_i x_{(i)}=\sum_i x_i$; hence $\sum_i|y_{(i)}-x_{(i)}|=\sum_i(y_{(i)}-x_{(i)})=y_j-x_j$. \qed

\emph{Claim (c).} If $x,y$ differ only in entry $j$, then $|\TM(x)-\TM(y)|\le|x_j-y_j|/(K-2b)$.

\emph{Proof.} $|\TM(x)-\TM(y)|\le\frac{1}{K-2b}\sum_{i=b+1}^{K-b}|x_{(i)}-y_{(i)}|\le\frac{1}{K-2b}\sum_{i=1}^{K}|x_{(i)}-y_{(i)}|$, and (b). \qed

\emph{Claim (d).} For any $x,y\in\R^K$, $\max_i|x_{(i)}-y_{(i)}|\le\|x-y\|_\infty$, and $|\TM(x)-\TM(y)|\le\|x-y\|_\infty$.

\emph{Proof.} Let $\epsilon=\|x-y\|_\infty$. Then $y\le x+\epsilon\mathbf 1$ and $x\le y+\epsilon\mathbf 1$ entrywise, so by (a) $y_{(i)}\le x_{(i)}+\epsilon$ and $x_{(i)}\le y_{(i)}+\epsilon$. Averaging over kept indices gives the second statement. \qed

\subsection*{Lemma 2 (sensitivity of a client's gap)}
\emph{Claim.} Fix the predictions $p(v)\in[0,1]$ of client $k$'s validation nodes. If the sensitive attribute of one validation node changes, $d_k$ changes by at most $2/n_{0,k}$.

\emph{Proof.} Write $n_0=n_{0,k}\ge1$. Let the node have prediction $p$, and consider the group it joins, with $N\ge0$ nodes and mass $S=\sum p(v)\in[0,N]$ before the change. Its rate changes from $S/\max(N,n_0)$ to $(S+p)/\max(N+1,n_0)$. Since $N$ and $n_0$ are integers, either $N+1\le n_0$ or $N\ge n_0$.
\begin{itemize}
\item If $N+1\le n_0$, both denominators equal $n_0$, and the change is $p/n_0\le1/n_0$.
\item If $N\ge n_0\ge1$, the change is $\frac{S+p}{N+1}-\frac{S}{N}=\frac{pN-S}{N(N+1)}$. Since $0\le pN\le N$ and $0\le S\le N$, $|pN-S|\le N$, and the change is at most $1/(N+1)\le1/n_0$ in absolute value.
\end{itemize}
The group the node leaves goes from $N'$ nodes with mass $S'$ to $N'-1$ nodes with mass $S'-p$; this is the computation above with the roles of the two states exchanged (apply it to the group of $N'-1$ nodes that the node joins), so the change is again at most $1/n_0$. The other group's rate and the predictions are unchanged. Hence $|\Delta(m_{k,1}-m_{k,0})|\le2/n_0$, and clipping to $[-1,1]$ is 1-Lipschitz. \qed

\subsection*{Theorem 3 (attribute-level differential privacy)}
\emph{Setting.} Training lasts $T$ rounds and statistics are released in the $R=\lceil T/P\rceil$ rounds $t$ with $t\bmod P=0$. Neighbouring datasets $D,D'$ differ only in the attribute of one validation node $v$ of one honest client $k$. We fix the training data of all clients, the initialisation and all randomness of local training; Byzantine clients' submissions are arbitrary (possibly randomised) functions of public information, which includes all released values, and of their own data. At least one server, say $A$, is honest. The released values are $\gt^{(j)}=\TM(\tilde d^{(j)})+\xi_A^{(j)}+\xi_B^{(j)}$, $j=1,\dots,R$, with independent $\xi_A^{(j)}\sim\mathcal{N}(0,\sigma_d^2)$, and the update sums $U^t$.

\emph{Step 1: the update channel does not depend on $s_v$ given the releases.} The model $\theta^{t}$ is determined by $\theta^{t-1}$, $\lambda^{t-1},\omega^{t-1}$ and the accepted updates. Each honest client's update depends on $\theta^{t-1}$, $(\lambda^{t-1},\omega^{t-1})$, its training data and its local randomness; the local objective uses training nodes only, and the backbone does not read $s$. Byzantine updates depend on public information. $(\lambda,\omega)$ are functions of the released gaps. By induction over rounds, every model, update and acceptance decision is the same function of the release history under $D$ and $D'$. In particular, given the history of the first $j-1$ releases, the inputs to the $j$-th release differ between $D$ and $D'$ only in $d_k$, by at most $2/n_{0,k}\le2/n_{0,\min}$ (Lemma 2), where we used that the predictions are identical. Honest client $k$ is accepted in both cases ($d_k\in[-1,1]$ and the norm check does not involve $s_v$), so the vectors $\tilde d^{(j)}$ differ in one entry.

\emph{Step 2: one release.} By Lemma 1(c), $|\TM(\tilde d^{(j)}(D))-\TM(\tilde d^{(j)}(D'))|\le\DTM=2/(n_{0,\min}(K-2b))$. Consider the adversary's view in the worst case that server $B$ is curious: it knows $\xi_B^{(j)}$ and therefore observes $\TM+\xi_A^{(j)}$. By the security of the two-party protocol against a semi-honest server, $B$'s view of the computation is simulatable from its own input and the opened values, so it learns nothing about $\tilde d$ beyond $\TM+\xi_A$ (and $U^t$, which by Step 1 adds nothing that depends on $s_v$ given the history). The map $D\mapsto\TM(\tilde d^{(j)}(D))+\xi_A^{(j)}$ is a Gaussian mechanism with sensitivity $\DTM$ and standard deviation $\sigma_d=z\DTM$, which is $(\alpha,\alpha\DTM^2/(2\sigma_d^2))=(\alpha,\alpha/(2z^2))$-RDP for every $\alpha>1$~[Mironov 2017]. If both servers are honest, the observer sees noise of variance $2\sigma_d^2$, and the guarantee is stronger; we use the weaker one.

\emph{Step 3: composition.} The $j$-th release is a Gaussian mechanism conditionally on the history of the earlier releases (Step 1), so adaptive composition of RDP over $R$ releases gives $(\alpha,R\alpha/(2z^2))$-RDP, and the conversion of Canonne, Kamath and Steinke gives $(\varepsilon,\delta)$-DP with $\varepsilon=\min_{\alpha>1}[R\alpha/(2z^2)+\log\frac{\alpha-1}{\alpha}-\frac{\log\delta+\log\alpha}{\alpha-1}]$. The code computes $z$ by bisection on this expression over a fixed grid of orders. The bias-corrected magnitude $\Dhat$, $\lambda$, $\omega$, the stored model and the returned model are functions of the releases and of quantities that by Step 1 do not depend on $s_v$ given the releases, hence post-processing. \qed

\emph{What is not covered.} The proof treats the noise as continuous and the trimmed mean as exact. A deployment over $\mathbb{Z}_{2^{64}}$ computes the trimmed mean of fixed-point values and must sample discrete noise; the sensitivity then increases by the rounding error and the RDP of a sum of two discrete Gaussians must be bounded separately (Canonne et al.; Kairouz et al.). The proof also assumes a semi-honest server.

\subsection*{Theorem 4 (update influence)}
Let $A$ be the accepted set, $B\subseteq A$ the accepted Byzantine clients with $|B|=f'$, $H=A\setminus B$, and suppose every honest client is accepted, so $|H|\ge K-f$. With $\bar\Delta_H,\bar\Delta_B$ the means over $H$ and $B$, $U/|A|=\frac{|H|}{|A|}\bar\Delta_H+\frac{f'}{|A|}\bar\Delta_B$, so $U/|A|-\bar\Delta_H=\frac{f'}{|A|}(\bar\Delta_B-\bar\Delta_H)$. Every accepted update has norm at most $C$, so $\|\bar\Delta_B-\bar\Delta_H\|_2\le2C$, and $|A|\ge|H|\ge K-f$, $f'\le f$. \qed

\subsection*{Theorem 5 (falsified gaps)}
Let $f\le b$ Byzantine clients submit arbitrary values; values outside $[-1,1]$ are replaced by $0$, so every entry of $\tilde d$ lies in $[-1,1]$, and the $K-f$ honest entries are the honest gaps $h_{(1)}\le\dots\le h_{(K-f)}$.
Fix $i$ with $b+1\le i\le K-b$.
\begin{itemize}
\item Lower bound. The entries $\tilde d_{(1)},\dots,\tilde d_{(i)}$ are $i$ entries $\le\tilde d_{(i)}$; at most $f$ of them are Byzantine, so at least $i-f$ honest gaps are $\le\tilde d_{(i)}$. Since $i-f\ge b+1-f\ge1$, the $(i-f)$-th smallest honest gap satisfies $h_{(i-f)}\le\tilde d_{(i)}$.
\item Upper bound. The entries $\tilde d_{(i)},\dots,\tilde d_{(K)}$ are $K-i+1$ entries $\ge\tilde d_{(i)}$; at least $K-i+1-f$ of them are honest. Since $i\le K-b\le K-f$, this number is at least $1$, and the $(K-i+1-f)$-th largest honest gap, which is $h_{(i)}$, satisfies $h_{(i)}\ge\tilde d_{(i)}$.
\end{itemize}
Averaging $h_{(i-f)}\le\tilde d_{(i)}\le h_{(i)}$ over $i=b+1,\dots,K-b$ gives the two inner inequalities of Eq.~(6) of the main text; the outer ones follow from $h_{(i-f)}\ge h_{(b+1-f)}$ and $h_{(i)}\le h_{(K-b)}$ for kept $i$. \qed

\emph{Consequences.} (i) is immediate. (ii) If all honest gaps except at most $b-f$ exceed $\tau'$, then $h_{(b-f+1)}>\tau'$ and so $\TM(\tilde d)>\tau'$; on the event $|\xi_A+\xi_B|\le r$ the released $\gt$ exceeds $\tau'-r$. (iii) $\TM(\tilde d)$ and the analogous honest-only quantity both lie in $[h_{(b-f+1)},h_{(K-b)}]$.

\emph{Breakdown.} If $f=b+1$ and all liars report $-1$, the smallest $b+1$ entries are $-1$ (when no honest gap equals $-1$), one of them is kept, and $\TM$ is the average of $-1$ and $K-2b-1$ honest values, which can lie outside the honest range. Each kept liar value can shift $\TM$ by up to $2/(K-2b)$ relative to an honest value.

\subsection*{Theorem 6 (returned model)}
For $\zeta\sim\mathcal{N}(0,s^2)$, $\Pr[|\zeta|>r]\le2\exp(-r^2/(2s^2))$. With $s^2=2\sigma_d^2$ and $r=r_R=2\sigma_d\sqrt{\ln(2R/\beta)}$ the right-hand side is $2\exp(-\ln(2R/\beta))=\beta/R$. A union bound over the $R$ independent releases gives $\Pr[\forall j:|\gt^{(j)}-\TM(\tilde d^{(j)})|\le r_R]\ge1-\beta$. A stored model satisfies $\Dhat\le\tau$, which by the definition $\Dhat=\sqrt{\max(0,\gt^2-2\sigma_d^2)}$ is equivalent to $\gt^2\le\tau^2+2\sigma_d^2$, i.e.\ $|\gt|\le\tau_{\mathrm{eff}}$. On the event, $|\TM(\tilde d(\hat\theta))|\le|\gt|+r_R\le\tau_{\mathrm{eff}}+r_R$. If $f\le b$, Theorem 5 gives $h_{(b-f+1)}\le\TM\le h_{(K-b)}$, hence $h_{(b-f+1)}(\hat\theta)\le\tau_{\mathrm{eff}}+r_R$ and $h_{(K-b)}(\hat\theta)\ge-(\tau_{\mathrm{eff}}+r_R)$. Note that the event concerns the noise only, so the bound holds whatever the liars do. \qed

\subsection*{Proposition 7 (validation to population)}
Under the stated assumption, conditionally on the group sizes $n_{k,a}\ge n_{0,k}$, $m_{k,a}$ is the mean of $n_{k,a}$ independent values in $[0,1]$ with mean $\mu_{k,a}$, so Hoeffding's inequality gives $\Pr[|m_{k,a}-\mu_{k,a}|>e]\le2\exp(-2n_{k,a}e^2)$. Setting this to $\beta'/(2K)$ gives $e_{k,a}=\sqrt{\ln(4K/\beta')/(2n_{k,a})}$, and a union bound over the $2K$ client-groups gives the event with probability $\ge1-\beta'$. On it, each $|d_k-(\mu_{k,1}-\mu_{k,0})|\le2\max e_{k,a}$ (the clip does not bind and is 1-Lipschitz), and Lemma 1(d) transfers this to the trimmed means. \qed

The assumption that $\hat\theta$ is independent of the validation draws fails for \rt{}, as discussed in the main text.

\subsection*{Differencing attack (closed form)}
For one release, a test between $\mathcal{N}(\mu,s^2)$ and $\mathcal{N}(\mu+\Delta,s^2)$ with equal priors is optimal when it thresholds at $\mu+\Delta/2$, and its balanced accuracy is $\Phi(|\Delta|/(2s))$. With both servers honest $s=\sqrt2\sigma_d$; with one colluding server $s=\sigma_d$. If the target's contribution is unchanged over $R$ releases and the attacker knows everything else, the sum of the $R$ observations is sufficient, with mean shift $R\Delta$ and standard deviation $\sqrt R s$, which gives $\Phi(\sqrt R|\Delta|/(2s))$. With $|\Delta|\le\DTM$ and $\sigma_d=z\DTM$ these values depend only on $z$.

\clearpage
\section{Additional tables}

\begin{table}[tbp]
\centering
\scriptsize
\setlength{\tabcolsep}{3pt}
\caption{Falsified gap reports against CLAMP-RT: test AUC, mean$\pm$sd over 3 seeds (companion of the DPD table in the main text). $\ast$: $f>b$.}
\label{tab:s_e4auc}
\begin{adjustbox}{max width=\textwidth}
\begin{tabular}{llcccccc}
\toprule
Dataset & $K$ & $f$ & Honest & hide & inflate & inlier & zero \\
\midrule
Bail & 5 & 1 & 0.781$\pm$0.046 & 0.781$\pm$0.047 & 0.781$\pm$0.046 & 0.781$\pm$0.045 & 0.781$\pm$0.046 \\
 &  & 2$^\ast$ &  & 0.764$\pm$0.038 & 0.775$\pm$0.038 & 0.781$\pm$0.046 & 0.781$\pm$0.046 \\
 & 10 & 1 & 0.743$\pm$0.019 & 0.743$\pm$0.019 & 0.743$\pm$0.019 & 0.743$\pm$0.019 & 0.743$\pm$0.020 \\
 &  & 2 &  & 0.743$\pm$0.019 & 0.743$\pm$0.019 & 0.743$\pm$0.019 & 0.743$\pm$0.019 \\
 &  & 3$^\ast$ &  & 0.732$\pm$0.020 & 0.743$\pm$0.017 & 0.744$\pm$0.019 & 0.744$\pm$0.019 \\
\midrule
Credit & 5 & 1 & 0.764$\pm$0.005 & 0.766$\pm$0.004 & 0.765$\pm$0.005 & 0.766$\pm$0.004 & 0.766$\pm$0.004 \\
 &  & 2$^\ast$ &  & 0.734$\pm$0.004 & 0.629$\pm$0.042 & 0.767$\pm$0.004 & 0.767$\pm$0.005 \\
 & 10 & 1 & 0.759$\pm$0.006 & 0.760$\pm$0.006 & 0.759$\pm$0.006 & 0.760$\pm$0.006 & 0.760$\pm$0.005 \\
 &  & 2 &  & 0.760$\pm$0.005 & 0.759$\pm$0.005 & 0.760$\pm$0.005 & 0.760$\pm$0.005 \\
 &  & 3$^\ast$ &  & 0.760$\pm$0.008 & 0.668$\pm$0.035 & 0.760$\pm$0.005 & 0.761$\pm$0.005 \\
\midrule
Pokec-z & 5 & 1 & 0.755$\pm$0.025 & 0.755$\pm$0.024 & 0.751$\pm$0.025 & 0.751$\pm$0.027 & 0.748$\pm$0.031 \\
 &  & 2$^\ast$ &  & 0.685$\pm$0.011 & 0.730$\pm$0.043 & 0.757$\pm$0.029 & 0.752$\pm$0.035 \\
 & 10 & 1 & 0.768$\pm$0.002 & 0.768$\pm$0.003 & 0.768$\pm$0.003 & 0.768$\pm$0.003 & 0.768$\pm$0.004 \\
 &  & 2 &  & 0.768$\pm$0.003 & 0.769$\pm$0.006 & 0.767$\pm$0.003 & 0.768$\pm$0.003 \\
 &  & 3$^\ast$ &  & 0.751$\pm$0.007 & 0.761$\pm$0.002 & 0.767$\pm$0.002 & 0.768$\pm$0.003 \\
\midrule
Pokec-n & 5 & 1 & 0.771$\pm$0.029 & 0.771$\pm$0.029 & 0.774$\pm$0.030 & 0.770$\pm$0.026 & 0.773$\pm$0.029 \\
 &  & 2$^\ast$ &  & 0.706$\pm$0.005 & 0.730$\pm$0.011 & 0.770$\pm$0.030 & 0.773$\pm$0.035 \\
 & 10 & 1 & 0.768$\pm$0.003 & 0.766$\pm$0.004 & 0.764$\pm$0.003 & 0.765$\pm$0.001 & 0.765$\pm$0.004 \\
 &  & 2 &  & 0.767$\pm$0.002 & 0.767$\pm$0.006 & 0.765$\pm$0.004 & 0.766$\pm$0.003 \\
 &  & 3$^\ast$ &  & 0.755$\pm$0.001 & 0.757$\pm$0.006 & 0.765$\pm$0.003 & 0.766$\pm$0.004 \\
\bottomrule
\end{tabular}
\end{adjustbox}
\end{table}

\begin{table}[tbp]
\centering
\scriptsize
\setlength{\tabcolsep}{3pt}
\caption{CLAMP-RT under update attacks ($K=10$, 3 seeds): test DPD / AUC, mean$\pm$sd.}
\label{tab:s_e3}
\begin{adjustbox}{max width=\textwidth}
\begin{tabular}{llcccc}
\toprule
Dataset & Attack & $f=1$ & $f=2$ & $f=3$ & $f=4$ \\
\midrule
Bail & none ($f=0$) & \multicolumn{4}{l}{DPD 0.009$\pm$0.002, AUC 0.743$\pm$0.019} \\
 & ALIE & 0.010$\pm$0.004 / 0.74 & 0.007$\pm$0.002 / 0.72 & 0.006$\pm$0.001 / 0.72 & 0.006$\pm$0.001 / 0.72 \\
 & Fair.\ poison & 0.010$\pm$0.002 / 0.74 & 0.009$\pm$0.008 / 0.72 & 0.015$\pm$0.009 / 0.71 & 0.023$\pm$0.009 / 0.71 \\
 & IPM & 0.009$\pm$0.002 / 0.73 & 0.006$\pm$0.002 / 0.71 & 0.005$\pm$0.002 / 0.71 & 0.004$\pm$0.001 / 0.71 \\
 & Scaling & 0.009$\pm$0.002 / 0.74 & 0.009$\pm$0.002 / 0.74 & 0.008$\pm$0.002 / 0.74 & 0.009$\pm$0.002 / 0.74 \\
 & Sign flip & 0.009$\pm$0.003 / 0.73 & 0.009$\pm$0.009 / 0.69 & 0.009$\pm$0.009 / 0.68 & 0.012$\pm$0.010 / 0.66 \\
\midrule
Pokec-z & none ($f=0$) & \multicolumn{4}{l}{DPD 0.046$\pm$0.008, AUC 0.768$\pm$0.003} \\
 & ALIE & 0.045$\pm$0.008 / 0.77 & 0.042$\pm$0.011 / 0.76 & 0.036$\pm$0.016 / 0.75 & 0.038$\pm$0.010 / 0.75 \\
 & Fair.\ poison & 0.059$\pm$0.003 / 0.77 & 0.118$\pm$0.013 / 0.75 & 0.151$\pm$0.023 / 0.74 & 0.106$\pm$0.097 / 0.73 \\
 & IPM & 0.045$\pm$0.009 / 0.77 & 0.046$\pm$0.010 / 0.78 & 0.044$\pm$0.008 / 0.78 & 0.025$\pm$0.005 / 0.76 \\
 & Scaling & 0.046$\pm$0.008 / 0.77 & 0.046$\pm$0.008 / 0.77 & 0.045$\pm$0.008 / 0.77 & 0.046$\pm$0.008 / 0.77 \\
 & Sign flip & 0.043$\pm$0.008 / 0.76 & 0.032$\pm$0.015 / 0.75 & 0.025$\pm$0.015 / 0.74 & 0.017$\pm$0.012 / 0.71 \\
\bottomrule
\end{tabular}
\end{adjustbox}
\end{table}

\begin{table}[tbp]
\centering
\footnotesize
\caption{Width $h_{(K-b)}-h_{(b-f+1)}$ of the interval of Theorem 5 of the main text in the attacked release rounds with $f\le b$ (pre-noise honest gaps): mean and maximum.}
\label{tab:s_spread}
\begin{adjustbox}{max width=\textwidth}
\begin{tabular}{lccc}
\toprule
Dataset & $K$ & mean & max \\
\midrule
Bail & 5 & 0.028 & 0.057 \\
Bail & 10 & 0.050 & 0.084 \\
Credit & 5 & 0.055 & 0.085 \\
Credit & 10 & 0.047 & 0.071 \\
Pokec-z & 5 & 0.044 & 0.072 \\
Pokec-z & 10 & 0.053 & 0.077 \\
Pokec-n & 5 & 0.049 & 0.079 \\
Pokec-n & 10 & 0.048 & 0.090 \\
\bottomrule
\end{tabular}
\end{adjustbox}
\end{table}

\begin{table}[tbp]
\centering
\scriptsize
\setlength{\tabcolsep}{3pt}
\caption{Privacy sweep: test AUC, mean$\pm$sd over 3 seeds.}
\label{tab:s_e5auc}
\begin{adjustbox}{max width=\textwidth}
\begin{tabular}{cccccccccc}
\toprule
 & \multicolumn{3}{c}{Credit} & \multicolumn{3}{c}{Pokec-z} & \multicolumn{3}{c}{Pokec-n} \\
$\varepsilon$ & $P=1$ & $P=5$ & $P=10$ & $P=1$ & $P=5$ & $P=10$ & $P=1$ & $P=5$ & $P=10$ \\
\midrule
1 & 0.757$\pm$0.011 & 0.765$\pm$0.004 & 0.766$\pm$0.006 & 0.747$\pm$0.019 & 0.749$\pm$0.025 & 0.751$\pm$0.032 & 0.770$\pm$0.026 & 0.773$\pm$0.030 & 0.774$\pm$0.029 \\
2 & 0.757$\pm$0.011 & 0.766$\pm$0.004 & 0.767$\pm$0.006 & 0.753$\pm$0.018 & 0.749$\pm$0.031 & 0.756$\pm$0.029 & 0.770$\pm$0.022 & 0.766$\pm$0.026 & 0.775$\pm$0.031 \\
4 & 0.755$\pm$0.015 & 0.766$\pm$0.005 & 0.766$\pm$0.004 & 0.759$\pm$0.030 & 0.754$\pm$0.023 & 0.755$\pm$0.025 & 0.768$\pm$0.027 & 0.764$\pm$0.024 & 0.770$\pm$0.031 \\
8 & 0.763$\pm$0.007 & 0.766$\pm$0.005 & 0.765$\pm$0.005 & 0.752$\pm$0.026 & 0.753$\pm$0.020 & 0.749$\pm$0.030 & 0.767$\pm$0.029 & 0.771$\pm$0.032 & 0.774$\pm$0.033 \\
\bottomrule
\end{tabular}
\end{adjustbox}
\end{table}


\end{document}
